% GENERATED FILE. Edit canonical lessons, then run npm run generate:books.
% canonical-source-hash: fnv1a64:d670a45878a0bdb2
% canonical-lessons: PT-C472-esperto, PT-C472-astuto, PT-C472-ingenuo, PT-C472-desconfiar, PT-C472-desconfiado

\chapter{Espertos e ingénuos}
\label{ch:pt-espertos-e-ingenuos}
\hlchaptermodality{472}

\begin{chapteropening}
\textbf{By the end of this chapter:} \emph{I can describe clever, cunning and naive people.}

Complete the last lesson of chapter 472: I can describe clever, cunning and naive people.
\end{chapteropening}

\section[esperto]{esperto — clever, smart}

\label{lesson:PT-C472-esperto}



\begin{quote}
Before the new one: say the Portuguese for liar, then the Portuguese for to pretend.
\end{quote}

\subsection*{You'll want to know: esperto}
\textbf{esperto} — \textquotedblleft{}clever, smart\textquotedblright{}.

\begin{quote}
A minha filha é muito esperta e aprende tudo depressa.
\end{quote}

\begin{quote}
\emph{My daughter is very clever and learns everything quickly.}
\end{quote}

Quick-witted rather than learned; \emph{inteligente} is broader.

\begin{cousinweb}
\emph{esperto} goes back to Old Portuguese \emph{esperto} (EtymDB); it is usually traced to Latin \emph{expertus} \textquotedblleft{}tried, experienced\textquotedblright{}, but the origin is not fully settled.
\end{cousinweb}

\begin{grammarlens}[title={What you've built}]
One more word for telling, explaining and describing at B1.
\end{grammarlens}

\subsection*{Your turn}
\begin{itemize}
\raggedright
  \item \emph{Say it:} \emph{esperto}
  \item \emph{Say it:} \emph{esperto}, once more
  \item \emph{Say it:} a sentence of your own with \emph{esperto}
  \item \emph{Recall:} say the Portuguese for liar, then the Portuguese for to pretend, then say \emph{esperto} again
\end{itemize}

\begin{tcolorbox}[breakable,colback=teal!4,colframe=teal!35!black,title={Before you move on}]
What does esperto mean? (\textquotedblleft{}Clever, smart\textquotedblright{}.) Say it once more.
\end{tcolorbox}
\section[astuto]{astuto — cunning, shrewd}

\label{lesson:PT-C472-astuto}



\begin{quote}
Before the new one: say the Portuguese for to pretend, then the Portuguese for clever.
\end{quote}

\subsection*{You'll want to know: astuto}
\textbf{astuto} — \textquotedblleft{}cunning, shrewd\textquotedblright{}.

\begin{quote}
Nas histórias para crianças, a raposa é astuta.
\end{quote}

\begin{quote}
\emph{In children's stories, the fox is cunning.}
\end{quote}

Less informal than \emph{esperto}.

\begin{cousinweb}
\emph{astuto} comes from Latin \emph{astūtus} \textquotedblleft{}shrewd, crafty\textquotedblright{} (EtymDB), built on \emph{astus} \textquotedblleft{}craft, cunning\textquotedblright{}.
\end{cousinweb}

\begin{grammarlens}[title={What you've built}]
One more word for telling, explaining and describing at B1.
\end{grammarlens}

\subsection*{Your turn}
\begin{itemize}
\raggedright
  \item \emph{Say it:} \emph{astuto}
  \item \emph{Say it:} \emph{astuto}, once more
  \item \emph{Say it:} a sentence of your own with \emph{astuto}
  \item \emph{Recall:} say the Portuguese for to pretend, then the Portuguese for clever, then say \emph{astuto} again
\end{itemize}

\begin{tcolorbox}[breakable,colback=teal!4,colframe=teal!35!black,title={Before you move on}]
What does astuto mean? (\textquotedblleft{}Cunning, shrewd\textquotedblright{}.) Say it once more.
\end{tcolorbox}
\section[ingénuo]{ingénuo — naive}

\label{lesson:PT-C472-ingenuo}



\begin{quote}
Before the new one: say the Portuguese for clever, then the Portuguese for cunning.
\end{quote}

\subsection*{You'll want to know: ingénuo}
\textbf{ingénuo} — \textquotedblleft{}naive\textquotedblright{}.

\begin{quote}
Era tão ingénuo que acreditava em tudo.
\end{quote}

\begin{quote}
\emph{He was so naive that he believed everything.}
\end{quote}

In Brazil: \emph{ingênuo}. Not English \emph{ingenious}: it means naive.

\begin{cousinweb}
\emph{ingénuo} is a learned borrowing from Latin \emph{ingenuus} \textquotedblleft{}native, freeborn\textquotedblright{}, and so \textquotedblleft{}frank, open\textquotedblright{} (EtymDB; Lewis \& Short).
\end{cousinweb}

\begin{grammarlens}[title={What you've built}]
One more word for telling, explaining and describing at B1.
\end{grammarlens}

\subsection*{Your turn}
\begin{itemize}
\raggedright
  \item \emph{Say it:} \emph{ingénuo}
  \item \emph{Say it:} \emph{ingénuo}, once more
  \item \emph{Say it:} a sentence of your own with \emph{ingénuo}
  \item \emph{Recall:} say the Portuguese for clever, then the Portuguese for cunning, then say \emph{ingénuo} again
\end{itemize}

\begin{tcolorbox}[breakable,colback=teal!4,colframe=teal!35!black,title={Before you move on}]
What does ingénuo mean? (\textquotedblleft{}Naive\textquotedblright{}.) Say it once more.
\end{tcolorbox}
\section[desconfiar]{desconfiar — to suspect; to distrust}

\label{lesson:PT-C472-desconfiar}



\begin{quote}
Before the new one: say the Portuguese for cunning, then the Portuguese for naive.
\end{quote}

\subsection*{You'll want to know: desconfiar}
\textbf{desconfiar} — \textquotedblleft{}to suspect; to distrust\textquotedblright{}.

\begin{quote}
Desconfio que o meu filho não fez os trabalhos de casa.
\end{quote}

\begin{quote}
\emph{I suspect my son has not done his homework.}
\end{quote}

\emph{Desconfiar que} — to suspect that; \emph{desconfiar de alguém} — to distrust someone.

\begin{cousinweb}
\emph{desconfiar} is built from \emph{des-} and \emph{confiar} \textquotedblleft{}to trust\textquotedblright{}, which goes back to Latin \emph{cōnfīdere} \textquotedblleft{}to trust fully\textquotedblright{} (Lewis \& Short), built on \emph{fīdere} \textquotedblleft{}to trust\textquotedblright{}, a relative of \emph{fidēs} \textquotedblleft{}faith\textquotedblright{}.
\end{cousinweb}

\begin{grammarlens}[title={What you've built}]
One more word for telling, explaining and describing at B1.
\end{grammarlens}

\subsection*{Your turn}
\begin{itemize}
\raggedright
  \item \emph{Say it:} \emph{desconfiar}
  \item \emph{Say it:} \emph{desconfiar}, once more
  \item \emph{Say it:} a sentence of your own with \emph{desconfiar}
  \item \emph{Recall:} say the Portuguese for cunning, then the Portuguese for naive, then say \emph{desconfiar} again
\end{itemize}

\begin{tcolorbox}[breakable,colback=teal!4,colframe=teal!35!black,title={Before you move on}]
What does desconfiar mean? (\textquotedblleft{}To suspect; to distrust\textquotedblright{}.) Say it once more.
\end{tcolorbox}
\section[desconfiado]{desconfiado — suspicious, wary}

\label{lesson:PT-C472-desconfiado}



\begin{quote}
Before the new one: say the Portuguese for naive, then the Portuguese for to suspect.
\end{quote}

\subsection*{You'll want to know: desconfiado}
\textbf{desconfiado} — \textquotedblleft{}suspicious, wary\textquotedblright{}.

\begin{quote}
O cão fica desconfiado com pessoas novas.
\end{quote}

\begin{quote}
\emph{The dog gets wary of new people.}
\end{quote}

The participle of \emph{desconfiar}.

\begin{cousinweb}
\emph{desconfiado} is the participle of \emph{desconfiar} \textquotedblleft{}to distrust\textquotedblright{}, built on \emph{confiar}, from Latin \emph{cōnfīdere} \textquotedblleft{}to trust fully\textquotedblright{} (Lewis \& Short).
\end{cousinweb}

\begin{grammarlens}[title={What you've built}]
One more word for telling, explaining and describing at B1.
\end{grammarlens}

\subsection*{Your turn}
\begin{itemize}
\raggedright
  \item \emph{Say it:} \emph{desconfiado}
  \item \emph{Say it:} \emph{desconfiado}, once more
  \item \emph{Say it:} a sentence of your own with \emph{desconfiado}
  \item \emph{Recall:} say the Portuguese for naive, then the Portuguese for to suspect, then say \emph{desconfiado} again
\end{itemize}

\begin{tcolorbox}[breakable,colback=teal!4,colframe=teal!35!black,title={Before you move on}]
What does desconfiado mean? (\textquotedblleft{}Suspicious, wary\textquotedblright{}.) Say it once more.
\end{tcolorbox}
